% ==================== 2.4 利用 Diamond 软件学习对称元素 ====================
\subsection{利用 Diamond 软件学习对称元素}

2.2.2 节已从理论上介绍了晶体的对称操作与对称元素；本节进一步利用 Diamond 软件对
反演中心、旋转轴、旋转反演轴、镜面、螺旋轴与滑移面六类对称元素逐一建模仿真，以弄清
各对称元素的作用方式、等效位置的生成规律及其在空间群符号中的体现。
各空间群一般位置的等效坐标均参照《国际晶体学表》\cite{hahn2002itt}。

为使对称操作的效果清晰可辨，六组演示采用统一的建模策略：选取能够体现目标对称元素、
且对称操作尽可能简洁的空间群，晶胞内只输入一个位于一般位置（general position）的
Si 原子（分数坐标 $(0.2,\,0.3,\,0.4)$），由 Diamond 按照空间群的对称操作自动生成晶胞内
全部等效原子；通过单胞图与超晶胞图的对比观察原子排布规律，相应的对称操作便一目了然。
操作步骤与 2.3.1 节相同（新建文档 $\to$ 输入空间群与晶格常数 $\to$ 添加原子
$\to$ Build $\to$ Fill $\to$ Unit Cell / Range），观察时的注意事项如下：

\begin{enumerate}[label=（\arabic*）]
  \item 空间群与晶系、晶胞参数的取值必须自洽，晶胞参数按各小节参数表中的演示取值
    输入；
  \item 原子必须位于一般位置：若坐标落入特殊位置（如 $0$、$1/2$、$1/3$ 等），部分
    对称操作会将该原子映射到自身，晶胞内显示的原子数将少于一般位置的等效位置数，
    无法观察到完整的对称操作效果；
  \item 螺旋轴与滑移面含平移分量，只有在超晶胞中才能显示出``旋转（反映）$+$ 平移''
    的复合操作特征；反演中心、旋转轴与旋转反演轴的超晶胞图则用于观察对称操作在
    空间中的周期复现；
  \item 观察时可借助鼠标旋转视图、调整观察方向，并可利用 Add $\to$ Atom Labels
    （ATOM 编号）追踪对称操作前后原子的对应关系。
\end{enumerate}
